% Section 4.11, from lectures/L04/appendix/b_exercises.md.

\section{Exercises}
\label{c4:sec:exercises}

\begin{aside}
\textbf{How to work these.} Each exercise is tagged with the kinds of work it asks for:
\kinds{Paper} by hand, \kinds{CircuitVerse} in the browser, \kinds{VHDL} in a module of your own.

\textbf{How to check your work.} Every exercise that asks you to write a VHDL module ships a
self-checking testbench under \file{lectures/L04/exercises} in the companion repository. Write your
module in its \file{exercises/<module>} directory, using the entity name and the \textbf{port
order} the exercise specifies, then run it with GHDL; \secref{c2:app:c} has the three commands.

No FPGA board is needed for any exercise. Your job is to get the VHDL right, and the testbench is
how you confirm it. Short answers to the parts you work on paper are in
Appendix~\ref{app:answers}; reference solutions to the VHDL are published in the companion
repository, under \file{lectures/L04/solutions}. Try each exercise yourself before looking at
either.
\end{aside}

% ------------------------------------------------------------------------------------------
\exerciseset{Metastability Concepts}
\label{c4:set:1}

\exercise{Metastability in your own words}{Paper}
\label{c4:ex:1}
Explain metastability in your own words. In your explanation, describe:
\begin{itemize}
  \item What happens inside a D flip-flop when its input changes too close to the active clock
    edge.
  \item Why the flip-flop output may temporarily remain between a valid logic \code{0} and logic
    \code{1}.
  \item Why the output may take an unpredictable amount of time to settle.
  \item Why two downstream gates may interpret the unresolved output differently:
    \begin{itemize}
      \item One gate may interpret it as \code{0}.
      \item Another gate may interpret it as \code{1}.
    \end{itemize}
\end{itemize}
Explain why this disagreement can cause different parts of a digital circuit to enter inconsistent
states.

\exercise{The setup-and-hold window}{Paper}
\label{c4:ex:2}
Setup time and hold time define a window around the active clock edge during which the input of a
flip-flop must remain stable.

\noindent Explain:
\begin{itemize}
  \item Why a synchronously generated signal can normally be designed to satisfy these timing
    requirements.
  \item Why an asynchronous input, such as a mechanical pushbutton, has no fixed relationship to the
    system clock.
  \item Why an asynchronous input may therefore change during the setup-and-hold window regardless
    of how the surrounding synchronous logic is designed.
\end{itemize}
Conclude why the metastability risk from an asynchronous input must be managed rather than simply
avoided.

% ------------------------------------------------------------------------------------------
\exerciseset{Synchronizer Design}
\label{c4:set:2}

\exercise{Drawing a double-flop synchronizer}{Paper}
\label{c4:ex:3}
Draw a double-flop synchronizer for one asynchronous input.

\noindent The circuit must contain:
\begin{itemize}
  \item One asynchronous input.
  \item Two D flip-flops connected in series.
  \item A shared system clock.
  \item A synchronized output taken from the second flip-flop.
\end{itemize}
Explain in your own words:
\begin{itemize}
  \item Why the first flip-flop is exposed directly to the asynchronous input.
  \item Why the first flip-flop may become metastable.
  \item How the interval between the first and second clock edges gives the first stage time to
    settle.
  \item Why the output of the second flip-flop has a much lower probability of being metastable.
  \item Why the second stage, rather than the first, should be used by the rest of the design.
\end{itemize}
A synchronizer does not make the probability of metastability exactly zero. Explain why it
nevertheless reduces the risk to an acceptably low level in most designs.

\note[Tip] Consider how much resolution time the first flip-flop receives before its value is
sampled by the second flip-flop.

\exercise{A third synchronizer stage}{Paper}
\label{c4:ex:4}
A colleague suggests using three synchronizer flip-flops on every asynchronous input in a low-speed
design clocked at a few megahertz, just to be safe.

A third stage buys one more clock period for a metastable value to settle, and costs one more clock
cycle of latency on every synchronized input. At a few megahertz that extra clock period is an
enormous amount of time in metastability terms, which is exactly why the proposal sounds prudent,
and also why it may already be unnecessary.

\begin{parts}
  \item Work the trade out in both directions. Under what circumstances does it come out in favour
    of the third stage, and under what circumstances is the second stage already enough? Your
    answer has to turn on something other than the clock frequency, because frequency alone does
    not settle it: name what it does turn on.
  \item \Secref{c4:a4} got to MTBF from the resolution time available per stage. Without redoing
    that arithmetic, say which way each of these pushes the decision, and why:
    \begin{itemize}
      \item How often the asynchronous input actually changes.
      \item The metastability characteristics of the FPGA family you are targeting.
      \item What a synchronization failure would cost if one happened.
    \end{itemize}
  \item Give your recommendation for an ordinary low-speed design like the ones in this book. Then
    describe a design, real or imagined, where you would argue for the third stage instead, and say
    which of the factors in \textbf{b)} is doing the work in that case.
\end{parts}
\textbf{Check your answer covers} the two things a third stage provides, an additional clock period
of resolution time and a lower probability that an unresolved value reaches the functional logic,
and the two it costs, an additional cycle of latency and additional flip-flops. If your answer to
\textbf{a)} rests on the clock rate rather than on the required mean time between failures, read
\secref{c4:a4} again before moving on.

\exercise{Assert async, release sync}{Paper}
\label{c4:ex:5}
Explain the \textbf{assert asynchronously, release synchronously} reset pattern used in
\file{reset_sync.vhd} (\secref{c4:a5}).

\noindent Describe the two parts separately:
\begin{itemize}
  \item Asynchronous assertion:
    \begin{itemize}
      \item The reset takes effect immediately.
      \item It does not wait for a clock edge.
    \end{itemize}
  \item Synchronous release:
    \begin{itemize}
      \item The reset is removed only on an active clock edge.
      \item All affected logic leaves reset in a controlled relationship with the clock.
    \end{itemize}
\end{itemize}
Explain:
\begin{itemize}
  \item Why immediate reset assertion can be important when:
    \begin{itemize}
      \item The clock is stopped.
      \item The clock is unstable.
      \item The system must enter its reset state without delay.
    \end{itemize}
  \item Why releasing an asynchronous reset close to a clock edge can violate the recovery or
    removal timing requirements of a flip-flop.
  \item Why synchronizing the release prevents different flip-flops from leaving reset on different
    clock cycles.
\end{itemize}

% ------------------------------------------------------------------------------------------
\exerciseset{Building the Circuit}
\label{c4:set:3}

\exercise{\code{led_toggle_sync} by hand}{CircuitVerse, Paper}
\label{c4:ex:6}
Build the \code{led_toggle_sync} circuit from \secref{c4:a9} by hand in CircuitVerse.

\begin{parts}
  \item Add the external ports:
    \begin{itemize}
      \item Inputs:
        \begin{itemize}
          \item \code{clock}
          \item \code{reset_n}: asynchronous, active-low.
          \item \code{button_n}: active-low.
        \end{itemize}
      \item Output:
        \begin{itemize}
          \item \code{led}
        \end{itemize}
    \end{itemize}
  \item Build the two-flip-flop reset synchronizer:
    \begin{itemize}
      \item Connect the two flip-flops in series.
      \item Assert both stages asynchronously using the raw \code{reset_n} input.
      \item Allow the reset to propagate out of the chain synchronously.
      \item Use the synchronized reset output, \code{reset_s2_n}, throughout the rest of the
        circuit.
      \item Do not use the raw \code{reset_n} input directly outside the reset synchronizer.
    \end{itemize}
  \item Build the three-flip-flop button synchronizer and edge detector:
    \begin{itemize}
      \item Use the first two flip-flops as the double-flop synchronizer.
      \item Use the third flip-flop to store the previous synchronized button value.
      \item Compare the current and previous synchronized values.
      \item Detect a button press when:
        \begin{itemize}
          \item The current synchronized active-low button value is \code{0}.
          \item The previous synchronized value is \code{1}.
        \end{itemize}
      \item Use an AND gate with the current value inverted to generate \code{button_edge_s2}.
    \end{itemize}
    The falling-edge relationship for the active-low button is:
    \[ \bool{button\_edge\_s2} = \bool{previous} \cdot \bool{current}' \]
  \item Build the LED toggle logic:
    \begin{itemize}
      \item Add a flip-flop to store \code{led_s}.
      \item Reset \code{led_s} using \code{reset_s2_n}.
      \item Toggle \code{led_s} only when \code{button_edge_s2 = 1}.
      \item Connect the \code{led} output to \code{led_s}.
    \end{itemize}
  \item Simulate the complete design:
    \begin{itemize}
      \item Set the clock period to \code{1000 ms}.
      \item Assert and release \code{reset_n}.
      \item Press and release \code{button_n}.
      \item Observe the synchronized button stages.
      \item Observe \code{button_edge_s2}.
      \item Observe \code{led_s}.
    \end{itemize}
    Determine:
    \begin{itemize}
      \item Whether the LED toggles on the first clock edge after the physical button press.
      \item How many clock cycles are required for the button value to propagate through the
        synchronizer and edge detector.
      \item Whether the observed latency matches the behaviour predicted in \secref{c4:a9}.
    \end{itemize}
\end{parts}

% ------------------------------------------------------------------------------------------
\exerciseset{Debouncing and Timing}
\label{c4:set:4}

\exercise{Synchronized is not debounced}{Paper}
\label{c4:ex:7}
You have now built and watched the circuit, so this is a question about what you saw.

The three-flip-flop chain in Exercise~\ref{c4:ex:6} does three jobs at once:
\begin{itemize}
  \item Reduces the risk of metastability.
  \item Stores the previous synchronized button value.
  \item Detects a button press.
\end{itemize}
Explain why it does \textbf{not}, by itself, guarantee that the button is debounced.

\begin{parts}
  \item Distinguish the three things it is easy to run together:
    \begin{itemize}
      \item Synchronization: converts an asynchronous input into a signal suitable for use in the
        clock domain.
      \item Edge detection: generates a pulse when the sampled signal changes.
      \item Debouncing: prevents several physical transitions from being interpreted as several
        separate presses.
    \end{itemize}
    Which two of these does your circuit actually do?
  \item In your simulation, one press produced one clean toggle. That is a stronger claim than it
    looks: it holds only under a particular condition about where the bouncing fell relative to the
    clock edges. State that condition precisely, in terms of what your circuit sampled and what it
    therefore saw.

    Then explain why that condition is an accident of your clock period rather than a property of
    your circuit. Be specific about which of the two, the bounce or the sampling, your design
    actually controls.
  \item Your simulation ran at a \code{1000 ms} clock period; the FPGA runs at \code{50 MHz}. A push
    button's contacts typically bounce for something on the order of 1\,ms. Work out the
    consequences yourself:
    \begin{itemize}
      \item How many times does a \code{1000 ms} clock sample the button during that 1\,ms of
        bouncing?
      \item How many times does a \code{50 MHz} clock sample it?
      \item Each sampled transition looks like a fresh edge to your edge detector. Given those two
        numbers, predict what the LED does on the board for one physical press, and say why the
        identical circuit behaved impeccably in CircuitVerse.
    \end{itemize}
  \item Describe one fix of each kind:
    \begin{itemize}
      \item A hardware fix: for example, an RC low-pass filter followed by a Schmitt-trigger input.
      \item A purely digital fix:
        \begin{itemize}
          \item for example, requiring the synchronized signal to hold a new level for some minimum
            number of clock cycles before accepting it.
          \item or a counter that ignores further transitions for a fixed interval after a press,
            which is what Chapter~\ref{c7:ch}'s timer would give you.
        \end{itemize}
    \end{itemize}
  \item Explain why the double-flop synchronizer must still be kept even once a debounce method is
    added. What problem does debouncing not solve?
\end{parts}

% ------------------------------------------------------------------------------------------
\exerciseset{The Synchronizers as Modules of Their Own}
\label{c4:set:5}
The worked example puts each synchronizer in a module of its own, so every later design
instantiates the same two files instead of re-deriving them. That is why Chapter~\ref{c7:ch}'s
\code{walking_led} and Chapter~\ref{c8:ch}'s \code{fsm_led} both ask you to copy these two in, and
why Chapter~\ref{c8:ch}'s \code{seq_detect_mealy}, which has no button, takes the reset synchronizer
alone and leaves the other behind (\secref{c4:a9}).

The next two exercises are you writing those modules. \Secref{c4:a5} and \secref{c4:a6} give you
the logic; what they do not give you is the entity around it, and, for the second one, the step
from a single button to a vector of them. Write each module yourself before opening the worked
example's copy, then compare. Exercise~\ref{c4:ex:10} then composes both of them into a working
design.

\exercise{The reset synchronizer}{VHDL, Paper}
\label{c4:ex:8}
Write the reset synchronizer as a module named \code{reset_sync}.

\noindent The entity has:

\begin{booktable}{@{}p{6em}p{4.9em}p{5.2em}L@{}}
  \thead{Port} & \thead{Direction} & \thead{Type} & \thead{Description} \\
  \midrule
  \code{clock} & in & \code{std_logic} & 50\,MHz system clock. \\
  \code{reset_n} & in & \code{std_logic} & Active-low asynchronous reset, straight from the
    outside world. \\
  \code{reset_s2_n} & out & \code{std_logic} & Active-low synchronized reset, safe for the rest of
    the design to use. Goes low the instant \code{reset_n} does, with no clock edge involved, and
    returns high two rising edges after \code{reset_n} does. \\
\end{booktable}

\begin{parts}
  \item Implement \textbf{assert asynchronously, release synchronously} (\secref{c4:a5}):
    \begin{itemize}
      \item One internal signal for the first stage, plus the output for the second.
      \item One process, sensitive to both \code{clock} and \code{reset_n}.
      \item On \code{reset_n = '0'}, drive both stages low, outside the clocked branch.
      \item On a rising clock edge, shift a constant \code{'1'} in through the two stages.
    \end{itemize}
  \item The value forced on reset and the value shifted in on the clock are opposites, \code{'0'}
    and \code{'1'}. Explain in one sentence why that is not a contradiction, in terms of what each
    branch is for.
  \item Verify the design with its testbench (see the aside at the start of this section). It checks
    both halves of the pattern separately: that pulling \code{reset_n} low takes \code{reset_s2_n}
    low with no clock edge in between, and that releasing \code{reset_n} leaves \code{reset_s2_n}
    low until two rising edges have passed.
  \item Answer, in prose:
    \begin{itemize}
      \item The output is named \code{reset_s2_n} rather than \code{reset_n_sync}. What does
        \code{s2} say, and what would have to change about this module for that name to become
        wrong?
      \item Every other module in this book takes \code{reset_s2_n} as its reset. This one takes
        the raw \code{reset_n}. Why must exactly one module in a design be the exception?
    \end{itemize}
\end{parts}
\note[Tip] the release is the interesting half. On reset the process does not care about the clock
at all, so both stages go low together; on the clock it does not care about \code{reset_n}, so the
\code{'1'} takes one edge per stage to reach the output. Two edges of latency on release is the
price of the whole pattern.

\modulefigure{reset_sync}

\selfcheck name your entity \code{reset_sync}, with ports \code{clock}, \code{reset_n} (in) and
\code{reset_s2_n} (out), declared in that order; its testbench is in \file{exercises/reset_sync}.

\exercise{The button synchronizer}{VHDL, Paper}
\label{c4:ex:9}
Write the button synchronizer and edge detector as a module named \code{button_sync}.

Exercise~\ref{c4:ex:10} handles two buttons by writing a two-bit synchronizer. This one is the same
circuit with the width left open: one generic decides how many buttons it serves, so
\code{led_toggle_sync} can use it for one button and \code{fsm_led} for two without either file
being edited. See \secref{c4:a8} for the syntax and for why this module has a generic where
\code{reset_sync} has none.

\noindent The entity has one generic:

\begin{booktable}{@{}p{4.5em}p{10.4em}p{3.8em}L@{}}
  \thead{Generic} & \thead{Type} & \thead{Default} & \thead{Description} \\
  \midrule
  \code{COUNT} & \code{natural range 1 to 3} & \code{1} & The number of buttons, and therefore the
    width of both vector ports. \\
\end{booktable}

\noindent and these ports:

\begin{booktable}{@{}p{7.4em}p{4.9em}>{\raggedright\arraybackslash}p{11.4em}L@{}}
  \thead{Port} & \thead{Direction} & \thead{Type} & \thead{Description} \\
  \midrule
  \code{clock} & in & \code{std_logic} & 50\,MHz system clock. \\
  \code{reset_s2_n} & in & \code{std_logic} & Active-low, \textbf{already synchronized} reset.
    This module does not synchronize its own reset; Exercise~\ref{c4:ex:8}'s module did that. \\
  \code{button_n} & in & \code{std_logic_vector(COUNT - 1 downto 0)} & Active-low asynchronous
    push buttons. \\
  \code{button_edge_s2} & out & \code{std_logic_vector(COUNT - 1 downto 0)} & A one-cycle pulse
    on each button's falling edge, two clock edges after the press. Each bit is independent of
    every other. \\
\end{booktable}

\begin{parts}
  \item Build the three-stage chain from \secref{c4:a6}, now over vectors rather than single bits:
    \begin{itemize}
      \item Three internal \code{std_logic_vector(COUNT - 1 downto 0)} signals, one per stage.
      \item Stages 1 and 2 are the double-flop synchronizer; stage 3 stores the previous value of
        stage 2.
      \item On reset, set all three stages to \code{(others => '1')}, not \code{(others => '0')}.
      \item Derive the output with a single concurrent assignment, and note that the same
        expression works unchanged on vectors as on bits.
    \end{itemize}
    \textbf{Keep this one in mind.} Those first two flip-flops are the whole of a synchronizer and
    nothing else, and Chapter~\ref{c5:ch}'s Exercise~\ref{c5:ex:4} has you pull them back out as a
    module of their own, generic in width and in the value they reset to. Writing them inline here
    first is what makes that exercise a consolidation rather than a new idea.
  \item Explain why the reset value is all-ones. All-ones means ``released'', because the buttons
    are active-low, so an all-zeros reset starts the chain out claiming every button is already
    held down. Trace what that costs:
    \begin{itemize}
      \item With the buttons released and a reset to all-zeros, does the module emit a spurious
        pulse as the released state propagates through the three stages? Work it through the
        output expression edge by edge before you answer, rather than guessing.
      \item Now the case that actually breaks. A user is holding a button down at the moment the
        reset is released. What does the module report with an all-ones reset, and what does it
        report with all-zeros? Which of the two is the behaviour you want, and why?
    \end{itemize}
  \item Verify the design with its testbench. It instantiates your module \textbf{twice}, once at
    \mbox{\code{COUNT = 1}} and once at \mbox{\code{COUNT = 2}}, and checks that a press produces
    exactly one single-cycle pulse, that holding the button produces no further pulses, that a
    release produces none at all, that pressing one button of a pair never pulses the other, and
    --- the case part \textbf{b)} asks you to reason about --- that a button already held down when
    the reset releases produces exactly one pulse with an all-ones reset and none at all with
    all-zeros.
  \item Answer, in prose:
    \begin{itemize}
      \item \code{COUNT} is constrained to \code{1 to 3} rather than left an unbounded
        \code{natural}. What does the constraint buy, given that nothing in the architecture
        depends on the upper bound?
      \item This module is deliberately not merged with Exercise~\ref{c4:ex:8}'s, even though a
        design that has a button always needs both. Give the design that justifies the split, and
        say what it would have had to do instead.
    \end{itemize}
\end{parts}
\note[Tip] write the architecture for \code{COUNT = 1} in your head first, then check that every
line you wrote is already correct for a vector. \code{and}, \code{not} and the aggregate
\code{(others => '1')} all work element-wise, so the generalization should cost you no extra code at
all. If it does, that is worth a second look.

\modulefigure{button_sync}

\selfcheck name your entity \code{button_sync}, with generic \code{COUNT}
(\code{natural range 1 to 3}, default \code{1}), inputs \code{clock}, \code{reset_s2_n},
\code{button_n} (\code{std_logic_vector(COUNT - 1 downto 0)}) and output \code{button_edge_s2}
(\code{std_logic_vector(COUNT - 1 downto 0)}), declared in that order; its testbench is in
\file{exercises/button_sync}.

% ------------------------------------------------------------------------------------------
\exerciseset{Composing Them Into a Design}
\label{c4:set:6}

\exercise{Two synchronized LED toggles}{VHDL, Paper}
\label{c4:ex:10}
Now make Chapter~\ref{c3:ch}'s \code{led_toggle} metastability-safe, in VHDL, as a new module named
\code{led_toggle_sync2}.

Chapter~\ref{c3:ch}'s \code{led_toggle} toggled two LEDs from two push buttons, but fed their raw,
asynchronous signals straight into the edge detector: logically correct, but unsafe on real
hardware (Chapter~\ref{c3:ch}'s own ``Looking ahead'' flagged exactly this). Here you rebuild it
with every asynchronous input guarded by this chapter's double-flop synchronizer.

\noindent The entity has:

\begin{booktable}{@{}p{5.2em}p{4.9em}p{14.4em}L@{}}
  \thead{Port} & \thead{Direction} & \thead{Type} & \thead{Description} \\
  \midrule
  \code{clock} & in & \code{std_logic} & 50\,MHz system clock. \\
  \code{reset_n} & in & \code{std_logic} & Active-low asynchronous reset. \\
  \code{button_n} & in & \code{std_logic_vector(1 downto 0)} & Two active-low asynchronous push
    buttons. \\
  \code{led} & out & \code{std_logic_vector(1 downto 0)} & Two LEDs; \code{led(i)} toggles once per
    press of \code{button_n(i)}. \\
\end{booktable}

\begin{parts}
  \item Synchronize every asynchronous input:
    \begin{itemize}
      \item Apply the two-flip-flop reset synchronizer from \secref{c4:a5} (assert asynchronously,
        release synchronously).
      \item Give the buttons a double-flop synchronizer plus a third stage for edge detection, as
        in \secref{c4:a6}, now over a \code{std_logic_vector(1 downto 0)}, so both buttons are
        synchronized in parallel.
      \item Never let the raw \code{reset_n} or \code{button_n} reach the toggle logic; use only
        the synchronized reset and the per-button edge pulses.
    \end{itemize}
    \note[Tip] you can inline the synchronizer, or split it out into a subcomponent taking a two-bit
    \code{button_n}, using the instantiation syntax from \secref{c2:a6}. Inlining is the shorter
    route, and then the build is a single file:
\begin{shell}
ghdl -a --std=93 led_toggle_sync2.vhd led_toggle_sync2_tb.vhd
\end{shell}
    If you split it out instead, copy in the \file{reset_sync.vhd} and \file{button_sync.vhd} you
    wrote for Exercises~\ref{c4:ex:8} and~\ref{c4:ex:9} (instantiating \code{button_sync} at
    \code{COUNT = 2}), and name them first on the \code{ghdl -a} line, ahead of your own module, so
    GHDL sees them before the file that instantiates them:
\begin{shell}
ghdl -a --std=93 reset_sync.vhd button_sync.vhd \
                 led_toggle_sync2.vhd led_toggle_sync2_tb.vhd
\end{shell}
    See \secref{c2:c5}. Neither module is shipped anywhere in the repository, because writing them
    is Exercises~\ref{c4:ex:8} and~\ref{c4:ex:9}. If yours are not working yet, go back to
    \secref{c4:a5}, which prints \code{reset_sync}'s process in full, and \secref{c4:a6}, which
    gives \code{button_sync}'s three-flop chain and its edge equation.
  \item Build the toggle logic:
    \begin{itemize}
      \item One LED-state flip-flop per LED.
      \item On the synchronized reset, clear both LEDs.
      \item Toggle \code{led(i)} only when button \code{i}'s synchronized falling-edge pulse fires.
      \item Confirm \code{led(i)} responds to \code{button_n(i)} alone, never to the other button.
    \end{itemize}
  \item Verify the design with its testbench (see the aside at the start of this section). The
    testbench presses each button in turn and checks that:
    \begin{itemize}
      \item Reset clears both LEDs.
      \item No press means no toggle.
      \item Pressing button 0 toggles \code{led(0)} \textbf{only}, once the synchronizer latency has
        elapsed.
      \item Holding that button down does not toggle it again.
      \item Pressing button 1 then toggles \code{led(1)} only, leaving \code{led(0)} where it was.
    \end{itemize}
    Pay attention to the \textbf{latency}: the testbench checks the \emph{exact} edge, not merely
    that the LED settles eventually. A press must reach the toggle logic through two synchronizer
    flip-flops, so \code{led} is still unchanged after the first and second rising edges and toggles
    on the third. Count the flip-flops a press travels through and satisfy yourself that the third
    edge is the right one, and that a defensive extra pipeline stage would now be caught rather than
    tolerated.
  \item Answer, in prose, by comparing this design with your Chapter~\ref{c3:ch} \code{led_toggle}:
    \begin{itemize}
      \item What does the synchronizer change about pressing a real, bouncing button?
      \item Why was the raw Chapter~\ref{c3:ch} version unsafe to wire to a physical button at
        \code{50 MHz}?
      \item Your testbench drives clean, perfectly-timed transitions. Which of the two problems
        from this chapter, metastability or contact bounce, can a testbench therefore never
        demonstrate to you?
    \end{itemize}
\end{parts}

\modulefigure{led_toggle_sync2}

\selfcheck name your entity \code{led_toggle_sync2}, with inputs \code{clock}, \code{reset_n},
\code{button_n} (\code{std_logic_vector(1 downto 0)}) and output \code{led}
(\code{std_logic_vector(1 downto 0)}), declared in that order; its testbench is in
\file{exercises/led_toggle_sync2}. It presses each button in turn and checks that \code{led(i)}
toggles on \code{button_n(i)}'s falling edge, that it does so on the right clock edge rather than
merely settling to the right value eventually, and that it never fires the wrong LED; see
\secref{c2:app:c}.
